% Kakeya sets over a fixed finite field: an entropy formula for the growth rate.
% By MV2 of Munim, Inc. Prepared with AI agents working in MV2's name. Text CC BY 4.0.
\documentclass[11pt]{amsart}
\usepackage{amsmath,amssymb,amsthm}
\usepackage[margin=1in]{geometry}

\newtheorem{theorem}{Theorem}
\newtheorem{lemma}[theorem]{Lemma}
\newtheorem{corollary}[theorem]{Corollary}
\theoremstyle{remark}
\newtheorem{remark}[theorem]{Remark}
\newtheorem{question}[theorem]{Question}

\newcommand{\F}{\mathbb{F}}
\newcommand{\Hb}{\mathrm{H}}

\title[Kakeya sets over a fixed finite field]{Kakeya sets over a fixed finite field:\\
a new upper bound and an entropy formula for the growth rate}
\author{MV2 (Munim, Inc.)}
\address{Munim, Inc.}
\email{munimversion2@gmail.com}
\date{October 15, 2026}

\begin{abstract}
Let $K(n,q)$ be the least size of a Kakeya set in $\F_q^n$ and
$\beta_q=\lim_{n\to\infty}K(n,q)^{1/n}$. Kopparty, Lev, Saraf and Sudan proved
$q^2/(2q-1)\le\beta_q$ and gave several upper bounds, the best for $q=3$ being
$3\cdot 2^{-2/3}\approx 1.8899$ by a random construction. We prove
$\beta_q\le q\,2^{-(q-1)/q}$ for every prime power $q$ with an explicit set. This is below every
upper bound we found in the literature for $q\ge 4$ (for example $2.3784$ against $2.8284$ at
$q=4$, and $4.3620$ against $6.3321$ at $q=8$), and matches $1.8899$ explicitly at $q=3$. We then
show that the averaging behind the construction is lossless: $\beta_q=\inf_k 2^{h_k/k}$, where
$h_k$ is the least average entropy $\frac1q\sum_{t\in\F_q}\Hb(P(X)+tX)$ over maps
$P:\F_q^k\to\F_q^k$ with $X$ uniform, and $h_1=\log_2 q-(q-1)/q$. Hence our bound equals $\beta_q$
if and only if $h_k=kh_1$ for all $k$, and a single finite map with smaller average entropy would
improve it. Finally, Kakeya sets invariant under permutations of the coordinates have size at least
$(q\,2^{-(q-1)/q})^{n}/(n+1)^q$, and over $\F_3$ they cannot reach $1.8899^n$.
\end{abstract}

\maketitle

\noindent\emph{Disclosure.} This paper, its proofs and the accompanying code were researched and
written with AI agents working in MV2's name. All finite statements are checked by the tests in
the repository \texttt{github.com/munimtechnologies/kakeya-small-fields}.

\section{Introduction}

A \emph{Kakeya set} in $\F_q^n$ is a set containing a line $\{a+td:t\in\F_q\}$ in every
direction $d\neq 0$. Dvir proved that Kakeya sets have size at least $c_n q^n$, and the method of
multiplicities and Bukh--Chao's work give the sharp density $2^{-(n-1)}$ as $q\to\infty$ with $n$
fixed. The regime of a \emph{fixed} field and growing dimension is different. Since products of
Kakeya sets are Kakeya sets, $K(m+n,q)\le K(m,q)K(n,q)$ and the limit
\[
\beta_q=\lim_{n\to\infty}K(n,q)^{1/n}=\inf_n K(n,q)^{1/n}
\]
exists. Kopparty, Lev, Saraf and Sudan (KLSS) \cite{KLSS} proved $\beta_q\ge q^2/(2q-1)$ and
upper bounds $(q+1)/2$ for odd $q$ (their Theorem 6) and $q\,2^{-2/q}$ (their Theorem 8, by a
random construction), and called the case $q=3$ most intriguing. Kyureghyan, M\"uller and Wang
\cite{KMW} improved the bounds for even $q$. For $q=2$, $\beta_2=\sqrt2$.

\begin{theorem}\label{thm:main}
For every prime power $q$, $\beta_q\le q\,2^{-(q-1)/q}$, with an explicit construction.
\end{theorem}

The base $q\,2^{-(q-1)/q}=\frac q2 2^{1/q}$ is within a factor $1+0.19/q+O(q^{-2})$ of the lower
bound $q^2/(2q-1)$.

The construction uses the lines of KLSS, $\{u^2+tu\}$ coordinatewise, but measures points by
the sizes of their composition classes on blocks of coordinates and shifts the parameter of the
line independently on each block. The $q$ points of a line then share the cost, and each pays the
average over $t$ of the entropy of $x\mapsto x^2+tx$ rather than the largest image size.

For a map $P:\F_q^k\to\F_q^k$ and $X$ uniform on $\F_q^k$ put
\[
h(P)=\frac1q\sum_{t\in\F_q}\Hb\bigl(P(X)+tX\bigr),\qquad h_k=\min_P h(P),
\]
with $\Hb$ the Shannon entropy in bits.

\begin{theorem}\label{thm:formula}
For every prime power $q$, $\beta_q=\lim_k 2^{h_k/k}=\inf_k 2^{h_k/k}$, and
$h_1=\log_2q-(q-1)/q$.
\end{theorem}

\begin{corollary}\label{cor:iff}
$\beta_q=q\,2^{-(q-1)/q}$ if and only if $h_k=kh_1$ for every $k\ge1$. Any map
$P:\F_q^k\to\F_q^k$ with $h(P)<kh_1$ gives $\beta_q\le 2^{h(P)/k}<q\,2^{-(q-1)/q}$.
\end{corollary}

\begin{theorem}\label{thm:sym}
\textup{(a)} If $K\subseteq\F_q^n$ is a Kakeya set invariant under all permutations of the
coordinates and $q\mid n$, then $|K|\ge (q\,2^{-(q-1)/q})^n/(n+1)^q$.
\textup{(b)} There is $\delta>0$ such that every permutation-invariant Kakeya set in $\F_3^n$,
$n$ even, has $|K|\ge 2^{(c_3+\delta)n}/(n+1)^3$, where $c_3=\log_2 3-2/3$ and
$2^{c_3}\approx1.8899$.
\end{theorem}

We searched the literature for bounds on $\beta_q$ with $q$ fixed and for formulas like
Theorem~\ref{thm:formula} and found none; we would be glad to learn of earlier proofs.

\section{The construction}

Write $c_q=\log_2q-(q-1)/q$. For $y\in\F_q^m$ let $N(y)$ be the number of vectors with the same
symbol counts as $y$. For integers $m,r\ge1$ and $n=rm$, cut $x\in\F_q^n$ into blocks
$x^{(1)},\dots,x^{(r)}$ of length $m$ and set
\[
w(x)=\sum_{j=1}^r\log_2 N(x^{(j)}),\qquad
K=\{x: w(x)\le T\},\qquad T=rmc_q+m\log_2 q\sqrt{r\ln q}.
\]
$K$ is a union of at most $\binom{m+q-1}{q-1}^r$ products of composition classes, each of size at
most $2^T$, so $|K|\le\binom{m+q-1}{q-1}^r2^T$, and with $m\asymp n^{1/3}$ this is
$(q2^{-(q-1)/q})^{n+O(n^{2/3}\log n)}$. It remains to show that $K$ is Kakeya.

Let $g_t(x)=x^2+tx$ on $\F_q$, applied coordinatewise to vectors.

\begin{lemma}\label{lem:entropy}
For every probability distribution $\tau$ on $\F_q$,
$\sum_{t\in\F_q}\Hb(g_t\#\tau)\le q\log_2q-(q-1)$, with equality at the uniform distribution.
\end{lemma}

\begin{proof}
The left side is concave in $\tau$. It is invariant under translating $\tau$: for odd $q$,
$g_t(x+b)=g_{t+2b}(x)+\mathrm{const}$ and $2b$ runs over $\F_q$; in characteristic $2$,
$g_t(x+b)=g_t(x)+\mathrm{const}$. Averaging the translates gives the uniform distribution, which is
therefore a maximum. There, for odd $q$, every $g_t$ has $(q-1)/2$ fibres of size $2$ and one of
size $1$, so $\Hb=\log_2q-(q-1)/q$; for even $q$, $g_0$ is a bijection and each $g_t$, $t\ne0$, is
exactly two-to-one. Both give $q\log_2q-(q-1)$.
\end{proof}

\begin{proof}[Proof of Theorem~\ref{thm:main}]
Let $d\ne0$ have blocks $u_j$, and let $a$ have blocks $u_j^2+s_ju_j$ for shifts
$s_j\in\F_q$. Block $j$ of $a+td$ is $g_{s_j+t}(u_j)$, so
$w(a+td)=\sum_jv_j(s_j+t)$ with $v_j(s)=\log_2N(g_s(u_j))$. If $u_j$ has symbol distribution
$\tau_j$ then $v_j(s)\le m\Hb(g_s\#\tau_j)$, so by Lemma~\ref{lem:entropy}
$\frac1q\sum_sv_j(s)\le mc_q$. Choose the $s_j$ independently and uniformly. For fixed $t$ the
terms $v_j(s_j+t)\in[0,m\log_2q]$ are independent with mean at most $mc_q$, so by Hoeffding's
inequality $\Pr[w(a+td)>T]\le q^{-2}$. A union bound over $t$ shows that some choice of shifts puts
the whole line in $K$.
\end{proof}

For $q=3$ the shifts can be chosen by an explicit greedy rule (using $\pm x^2$ to realise all six
orders of the three class sizes on each block), giving an explicit set of size
$(3/4^{1/3})^{n+O(\sqrt{n\log n})}$; see the repository.

\begin{lemma}[limit of one-coordinate witnesses]\label{lem:limit}
For every $f:\F_q\to\F_q$ and $X$ uniform on $\F_q$,
$\sum_t\Hb(f(X)+tX)\ge q\log_2q-(q-1)$.
\end{lemma}

\begin{proof}
Two points $x\ne y$ collide under $x\mapsto f(x)+tx$ for exactly one $t$, so the fibres of the
$q$ maps contain $\binom q2$ pairs in total. A fibre of size $s$ lowers the entropy from $\log_2q$
by $(s\log_2s)/q$ and uses $\binom s2$ pairs, and $s\log_2s\le 2\binom s2$. So the total reduction
is at most $q-1$.
\end{proof}

\section{The entropy formula}

\begin{proof}[Proof of Theorem~\ref{thm:formula}]
(i) \emph{Subadditivity.} If $P$ acts on $\F_q^k$ and $P'$ on $\F_q^l$, then $P\times P'$ acts on
$\F_q^{k+l}$ and, $X$ splitting into independent uniform parts, $h(P\times P')=h(P)+h(P')$. Hence
$h_{k+l}\le h_k+h_l$, and by Fekete's lemma $\lim_k h_k/k=\inf_k h_k/k$.

(ii) \emph{$h_k\le\log_2K(k,q)$.} Let $K_0$ be a Kakeya set in $\F_q^k$. For $x\ne0$ let $P(x)$
satisfy $P(x)+tx\in K_0$ for all $t$, and let $P(0)\in K_0$. Every $P(X)+tX$ takes values in
$K_0$, so $h(P)\le\log_2|K_0|$. Therefore $\inf_k2^{h_k/k}\le\inf_kK(k,q)^{1/k}=\beta_q$.

(iii) \emph{$\beta_q\le2^{h_k/k}$.} For a distribution $\tau$ on $\F_q^k$ let
$\varphi(\tau)=\min_P\frac1q\sum_t\Hb(P(Y)+tY)$ with $Y\sim\tau$. For fixed $P$ each term is
concave in $\tau$, so $\varphi$ is concave. If $Y'=Y+b$ and $P'(y)=P(y-b)$, then
$P'(Y')+tY'=P(Y)+tY+tb$, so $\varphi$ is translation invariant. Averaging the translates of $\tau$
gives the uniform distribution, so $\varphi(\tau)\le h_k$ for all $\tau$.

Now repeat Section 2 over the alphabet $\Sigma=\F_q^k$. Write a point of $\F_q^{kmr}$ as $r$ blocks
of $m$ letters from $\Sigma$, let $N(y)$ count the words with the same letter counts as $y$, and
let $K=\{x:\sum_j\log_2N(x^{(j)})\le rmh_k+mk\log_2q\sqrt{r\ln q}\}$. For a direction $d$ whose
block $j$ has letter distribution $\tau_j$, pick $P_j$ attaining $\varphi(\tau_j)$ and a shift
$s_j$, and give $a$ the letter $P_j(u)+s_ju$ wherever block $j$ of $d$ has the letter $u$. Block
$j$ of $a+td$ has letter distribution equal to the law of $P_j(Y)+(s_j+t)Y$, $Y\sim\tau_j$, so its
cost is at most $m$ times that entropy, whose mean over a uniform $s_j$ is at most
$m\varphi(\tau_j)\le mh_k$. Hoeffding's inequality with range $mk\log_2q$ and a union bound over
$t$ give a line in $K$, and
$|K|\le\binom{m+q^k-1}{q^k-1}^r2^{rmh_k+mk\log_2q\sqrt{r\ln q}}$. Letting $m\asymp n^{1/3}$ gives
$\beta_q\le2^{h_k/k}$.

(iv) Lemma~\ref{lem:limit} gives $h_1\ge c_q$, and $P(x)=x^2$ attains it
(Lemma~\ref{lem:entropy}).

Combining (ii) and (iii), $\beta_q\le\inf_k2^{h_k/k}\le\beta_q$.
\end{proof}

\begin{proof}[Proof of Corollary~\ref{cor:iff}]
By subadditivity $h_k\le kh_1$. If $h_k=kh_1$ for all $k$ then $\beta_q=2^{h_1}=q2^{-(q-1)/q}$.
If $h(P)<kh_1$ for some $P$ on $\F_q^k$ then $\beta_q\le2^{h(P)/k}<2^{h_1}$ by (iii).
\end{proof}

\begin{remark}
$P$ chooses one line $\{(t,P(x)+tx)\}$ in each of the $q^k$ non-vertical directions $(1,x)$ of
$\F_q^{k+1}$, and $h(P)$ is the mean over the $q$ slices $\{t\}\times\F_q^k$ of the entropy of the
point where a uniformly random one of these lines meets the slice. With
$S(P)=\sum_t\sum_{\text{fibres}}s\log_2s$, $h(P)=k\log_2q-S(P)/q^{k+1}$, and $h(P)<kh_1$ exactly
when $S(P)>k(q-1)q^k$. The lower bound $\beta_q\ge q^2/(2q-1)$ gives
$h_k\ge k\log_2(q^2/(2q-1))$. For $q=3$, $k=2$ an improvement needs $36<S(P)\le39.80$; the
coordinatewise squares give $S=36$.
\end{remark}

\begin{question}
Is $h_k=kh_1$ for all $k$, that is, is $\beta_q=q\,2^{-(q-1)/q}$? This holds for $q=2$. The
smallest open case is $q=3$, $k=2$: is there $P:\F_3^2\to\F_3^2$ with $S(P)>36$?
\end{question}

\section{Permutation-invariant sets}

\begin{proof}[Proof of Theorem~\ref{thm:sym}]
(a) Let $u$ be a direction in which each symbol occurs $n/q$ times and $a+tu$ a line in $K$. Let
$\pi$ be the empirical joint distribution of the pairs $(u_i,a_i)$; its first marginal is uniform,
and the symbol distribution of $a+tu$ is the law of $Y+tX$, $(X,Y)\sim\pi$. Since $K$ contains the
composition class of $a+tu$, $|K|\ge2^{n\Hb(Y+tX)}/(n+1)^q$. Writing
$\pi(x,y)=\kappa(y\mid x)/q$ as the mixture $\sum_fw_f\pi_f$ of the laws $\pi_f$ of $(X,f(X))$,
with $w_f=\prod_x\kappa(f(x)\mid x)$, concavity and Lemma~\ref{lem:limit} give
$\sum_t\Hb(Y+tX)\ge\sum_fw_f\sum_t\Hb(f(X)+tX)\ge q\log_2q-(q-1)$, so some $t$ has
$\Hb(Y+tX)\ge c_q$.

(b) Let $u$ have $n/2$ zeros and $n/2$ ones, and let the symbols of $a$ have distribution $k_0$
where $u$ is $0$ and $k_1$ where $u$ is $1$. Point $t$ has distribution
$M_t(z)=(k_0(z)+k_1(z-t))/2$. It suffices to show $\max_t\Hb(M_t)>h(2/3)=c_3$ for all $k_0,k_1$,
where $h$ is the binary entropy; compactness then gives $\delta$. Suppose $\Hb(M_t)\le h(2/3)$
for all $t$. A distribution on three points with largest mass $p<2/3$ has entropy above $h(2/3)$
(if $p<1/2$ its entropy exceeds $1$; otherwise it is at least $h(p)$). So there are $z_t$ with
$k_0(z_t)+k_1(z_t-t)\ge4/3$, and $A+B\ge4$ where $A=\sum_tk_0(z_t)$, $B=\sum_tk_1(z_t-t)$. Since
$z_s=z_t$ forces $z_s-s\ne z_t-t$: if all $z_t$ coincide then $A\le3$, $B=1$; if exactly two
coincide then $A,B\le2$; if all differ then $A=1$, $B\le3$. Hence $A+B=4$, each $M_t$ has largest
mass exactly $2/3$, and $\Hb(M_t)\le h(2/3)$ forces $M_t$ to have the shape $(2/3,1/3,0)$. In the
first case $k_0$ is a point mass and $k_1$ uniform, so $M_t=(2/3,1/6,1/6)$; the third case is
symmetric; in the second $k_0,k_1$ are point masses and all masses of $M_t$ lie in
$\{0,1/2,1\}$. Each case is a contradiction.
\end{proof}

The set of points of $\F_3^n$ using at most two symbols is a permutation-invariant Kakeya set of
size $3\cdot2^n-3$ (each $g_t(u)$ uses at most two symbols), so for such sets the best base over
$\F_3$ lies in $(1.8899,2]$; we conjecture it is $2$.

\begin{thebibliography}{9}
\bibitem{BC} B.~Bukh and T.-W.~Chao, Sharp density bounds on the finite field Kakeya problem,
\emph{Discrete Analysis} 2021:26, arXiv:2108.00074.
\bibitem{Csiszar} I.~Csisz\'ar, The method of types, \emph{IEEE Trans. Inform. Theory} 44 (1998),
2505--2523.
\bibitem{DKSS} Z.~Dvir, S.~Kopparty, S.~Saraf and M.~Sudan, Extensions to the method of
multiplicities, with applications to Kakeya sets and mergers, FOCS 2009, arXiv:0901.2529.
\bibitem{Dvir} Z.~Dvir, On the size of Kakeya sets in finite fields, \emph{J. Amer. Math. Soc.} 22
(2009), 1093--1097.
\bibitem{GGTW} B.~Georgiev, J.~G\'omez-Serrano, T.~Tao and A.~Z.~Wagner, Mathematical exploration
and discovery at scale, arXiv:2511.02864.
\bibitem{KLSS} S.~Kopparty, V.~F.~Lev, S.~Saraf and M.~Sudan, Kakeya-type sets in finite vector
spaces, \emph{J. Algebraic Combin.} 34 (2011), 337--355, arXiv:1003.3736.
\bibitem{KMW} G.~Kyureghyan, P.~M\"uller and Q.~Wang, On the size of Kakeya sets in finite vector
spaces, \emph{Electron. J. Combin.} 20(3) (2013), \#P36, arXiv:1302.5591.
\end{thebibliography}

\end{document}
